\chapter{Przestrzenie wektorowe}
\label{ch:lin-przestrzenie-wektorowe}
\index{przestrzen wektorowa@Przestrzeń wektorowa}

\noindent\textbf{Cel rozdziału.}
Algebra liniowa bada obiekty, w których można dodawać elementy i mnożyć je
przez skalary, a także odwzorowania zachowujące te działania.
Zwykłe strzałki na płaszczyźnie są ważnym przykładem, ale nie definicją:
wektorami mogą być wielomiany, macierze, funkcje i ciągi.
Chcemy zrozumieć, co we wszystkich tych przykładach oznaczają
\emph{kombinacja liniowa}, \emph{niezależność liniowa} i \emph{baza}.
Wyjaśnimy również, dlaczego w przestrzeniach nieskończenie wymiarowych
baza algebraiczna nie jest tym samym co rozwinięcie w nieskończony szereg.

W całym rozdziale \(K\) oznacza ustalone \emph{ciało przemienne}
z \(0_K\neq1_K\), a \(V\) przestrzeń wektorową nad \(K\).
Przykładami są \(\mathbb Q,\mathbb R,\mathbb C\), ale także ciała skończone,
na przykład \(\mathbb F_2=\{0,1\}\).
Nie zakładamy, że \(K\) jest uporządkowane ani że ma charakterystykę zero.
Jeżeli potrzebujemy konkretnego ciała, podajemy je wprost.
Definicje i dowody teorii algebraicznej nie wymagają
topologii ani pojęcia zbieżności.
Jedynie w jednym z przykładów odwołamy się do znanej
z analizy własności sum funkcji ciągłych.
Pojęcia ciała i grupy przemiennej są znane z algebry abstrakcyjnej;
dla samodzielności przytaczamy dokładnie te aksjomaty, które są tutaj potrzebne.

\section{Dlaczego przestrzenie wektorowe?}
\label{sec:lin-motywacja}

Najprostszy przykład stanowi \(K^n\), zbiór \(n\)-elementowych kolumn
o współczynnikach w \(K\). Dodawanie kolumn i mnożenie przez skalar
wykonujemy współrzędna po współrzędnej. Jednak podobne rachunki można
przeprowadzać na wielomianach:
\[
 (1+X)+(-1+2X^2)=X+2X^2,\qquad
 3(1+X)=3+3X .
\]
W pierwszym wzorze przyjmujemy \(K=\mathbb R\).
Przy funkcjach \(f,g\colon S\to K\) działania
\((f+g)(s)=f(s)+g(s)\) i \((af)(s)=af(s)\) mają tę samą
strukturę. Nie istnieje żaden powód, by ogólną teorię
ograniczać do strzałek lub do ustalonej liczby współrzędnych.

\begin{definicja}[Przestrzeń wektorowa]
\label{def:lin-przestrzen-wektorowa}
\index{przestrzen wektorowa!aksjomaty@Przestrzeń wektorowa!aksjomaty}
\emph{Przestrzenią wektorową nad \(K\)} nazywamy zbiór \(V\)
z działaniem \(+\colon V\times V\to V\) oraz
mnożeniem przez skalary \(K\times V\to V\), \((a,v)\mapsto av\),
spełniający następujące warunki:
\begin{enumerate}[label=(\roman*)]
\item \((V,+)\) jest grupą przemienną; jej element neutralny
      oznaczamy przez \(0_V\), a element przeciwny do \(v\) przez \(-v\);
\item \(a(u+v)=au+av\) dla wszystkich \(a\in K\) i \(u,v\in V\);
\item \((a+b)v=av+bv\) dla wszystkich \(a,b\in K\) i \(v\in V\);
\item \((ab)v=a(bv)\) dla wszystkich \(a,b\in K\) i \(v\in V\);
\item \(1_Kv=v\) dla każdego \(v\in V\).
\end{enumerate}
Elementy \(V\) nazywamy \emph{wektorami}, a elementy \(K\)
\emph{skalarami}. Piszemy zwykle \(0\) zamiast \(0_K\) lub
\(0_V\), gdy rozróżnienie nie jest potrzebne.
\end{definicja}

Przemienność dodawania jest częścią definicji. Gdyby w punkcie
(i) założyć jedynie strukturę grupy, nie otrzymalibyśmy
automatycznie standardowego pojęcia przestrzeni wektorowej.
Ciało jest tutaj istotne: dzielenie przez niezerowy skalar
będzie kluczowe w dowodzie twierdzenia o wymianie wektorów.

\begin{lemat}[Podstawowe konsekwencje aksjomatów]
\label{lem:lin-aksjomaty-konsekwencje}
Dla dowolnych \(a\in K\) i \(v\in V\) zachodzą:
\(0_Kv=0_V\), \(a0_V=0_V\), \((-a)v=-(av)=a(-v)\).
Ponadto, jeżeli \(av=0_V\), to \(a=0_K\) lub \(v=0_V\).
\end{lemat}
\begin{proof}
Z równości \((0_K+0_K)v=0_Kv+0_Kv\) i skracania w
grupie \((V,+)\) wynika \(0_Kv=0_V\).
Podobnie \(a(0_V+0_V)=a0_V+a0_V\), więc \(a0_V=0_V\).
Ponieważ \(av+(-a)v=(a-a)v=0_V\), wektor \((-a)v\)
jest przeciwny do \(av\). Także \(a(-v)+av=a(0_V)=0_V\).
Jeżeli \(a\neq0_K\) i \(av=0_V\), to po pomnożeniu
obu stron przez \(a^{-1}\) dostajemy
\(v=1_Kv=a^{-1}(av)=0_V\).
\end{proof}

\begin{uwaga}[Uwaga o charakterystyce]
\label{rem:lin-charakterystyka}
W ciele \(\mathbb F_2\) zachodzi \(-1=1\), zatem
\(-v=v\) dla każdego wektora nad tym ciałem.
Nie jest to sprzeczność: wówczas \(v+v=0\).
Z tego powodu w ogólnej algebrze liniowej nie wolno
bez dodatkowego założenia dzielić przez \(2\)
ani utożsamiać wektorów z geometrycznymi strzałkami
w przestrzeni euklidesowej.
\end{uwaga}

\section{Przykłady i kontrprzykłady}
\label{sec:lin-przyklady-przestrzeni}

\begin{przyklad}[Przestrzeń współrzędnych]
\label{prz:lin-kn}
Dla liczby całkowitej \(n\geq1\) zbiór \(K^n\) jest przestrzenią
wektorową z działaniami współrzędnymi:
\((x_i)_{i=1}^n+(y_i)_{i=1}^n=(x_i+y_i)_{i=1}^n\) oraz
\(a(x_i)_{i=1}^n=(ax_i)_{i=1}^n\).
Wektor zerowy ma wszystkie współrzędne równe \(0_K\).
Każdy z aksjomatów sprawdzamy współrzędna po współrzędnej;
na przykład \((a+b)(x_i)_i=((a+b)x_i)_i
=(ax_i+bx_i)_i=a(x_i)_i+b(x_i)_i\).
Przyjmujemy także \(K^0=\{0\}\), jednopunktową przestrzeń zerową.
\end{przyklad}

\begin{przyklad}[Wielomiany]
\label{prz:lin-wielomiany}
Niech \(K[X]\) będzie zbiorem wielomianów jednej zmiennej
o współczynnikach w \(K\). Zwykłe dodawanie i mnożenie
wielomianu przez skalar dają przestrzeń wektorową.
Dla \(d\geq0\) jej podzbiór
\(K[X]_{\leq d}:=\{0\}\cup
\{p\in K[X]\setminus\{0\}:\deg p\leq d\}\)
również jest zamknięty na te działania.
Umowa \(\deg0=-\infty\) pozwoliłaby zapisać ten podzbiór
bez dołączania \(\{0\}\).
Niezerowy wielomian stopnia \(d+1\) nie należy do
\(K[X]_{\leq d}\), choć należy do \(K[X]\).
W dalszej części wykażemy, że ta podprzestrzeń ma bazę
\(1,X,\ldots,X^d\), a cała \(K[X]\) ma bazę nieskończoną.
\end{przyklad}

\begin{przyklad}[Przestrzenie funkcji]
\label{prz:lin-funkcje}
Dla dowolnego zbioru \(S\) wszystkie funkcje \(S\to K\)
tworzą przestrzeń \(K^S\) z działaniami określonymi punktowo.
Dotyczy to również \(S=\varnothing\): istnieje wtedy
dokładnie jedna funkcja \(\varnothing\to K\), więc
\(K^\varnothing\) jest przestrzenią zerową.

Gdy \(K=\mathbb R\) i \(S=\mathbb R\), funkcje ciągłe
\(\mathbb R\to\mathbb R\) tworzą podprzestrzeń \(K^S\):
suma i skalarna wielokrotność funkcji ciągłych są ciągłe.
Natomiast zbiór funkcji \(f\colon\mathbb R\to\mathbb R\)
spełniających \(f(0)=1\) nie jest podprzestrzenią,
ponieważ nie zawiera funkcji zerowej.
\end{przyklad}

\begin{przyklad}[Zmiana ciała skalarów]
\label{prz:lin-c-nad-r}
Liczby zespolone tworzą przestrzeń wektorową zarówno
nad \(\mathbb C\), jak i nad \(\mathbb R\).
W pierwszym przypadku każdy \(z\in\mathbb C\)
ma postać \(z\cdot1\), natomiast w drugim
\(z=x+iy=x\cdot1+y\cdot i\), gdzie \(x,y\in\mathbb R\).
Co więcej, jeśli \(a+bi=0\) dla rzeczywistych \(a,b\),
to \(a=b=0\). W konsekwencji \((1)\) jest bazą nad \(\mathbb C\),
a \((1,i)\) jest bazą nad \(\mathbb R\).
Wymiar i niezależność liniowa zależą więc od wskazanego ciała.
\end{przyklad}

\begin{przyklad}[Wszystkie ciągi i ciągi o skończonym nośniku]
\label{prz:lin-ciagi}
Niech \(\mathbb N_0=\{0,1,2,\ldots\}\).
Przestrzeń \(K^{\mathbb N_0}\) składa się ze
\emph{wszystkich} ciągów \((x_0,x_1,\ldots)\).
Jej podprzestrzeń
\[
 K^{(\mathbb N_0)}
  :=\{x\in K^{\mathbb N_0}:
          x_n=0\ \text{dla wszystkich poza skończenie wieloma }n\}
\]
składa się z ciągów o \emph{skończonym nośniku}.
Dodawanie takich ciągów nie wyprowadza poza tę podprzestrzeń:
nośnik sumy zawiera się w sumie dwóch skończonych nośników.
Ciąg stały \((1,1,\ldots)\) należy do \(K^{\mathbb N_0}\),
lecz nie należy do \(K^{(\mathbb N_0)}\).
To rozróżnienie okaże się istotne przy bazach.
\end{przyklad}

\begin{uwaga}[Moduł a przestrzeń wektorowa]
\label{rem:lin-moduly}
Przestrzeń wektorowa nad \(K\) jest modułem nad pierścieniem \(K\).
Nie każdy moduł jest jednak przestrzenią wektorową:
na przykład grupa addytywna \(\mathbb Z/2\mathbb Z\)
jest modułem nad \(\mathbb Z\), lecz nie przestrzenią wektorową
nad \(\mathbb Z\), ponieważ \(\mathbb Z\) nie jest ciałem.
Istnienie odwrotności każdego niezerowego skalara
odróżnia późniejsze twierdzenia o bazach przestrzeni wektorowych
od ogólnej teorii modułów.
\end{uwaga}

\section{Podprzestrzenie wektorowe}
\label{sec:lin-podprzestrzenie}

\begin{definicja}[Podprzestrzeń]
\label{def:lin-podprzestrzen}
Podzbiór \(W\subseteq V\) nazywamy \emph{podprzestrzenią wektorową},
jeśli \(W\) z działaniami odziedziczonymi z \(V\)
jest przestrzenią wektorową nad tym samym ciałem \(K\).
\end{definicja}

\begin{twierdzenie}[Kryterium podprzestrzeni]
\label{thm:lin-kryterium-podprzestrzeni}
Dla podzbioru \(W\subseteq V\) następujące warunki są równoważne:
\begin{enumerate}[label=(\roman*)]
\item \(W\) jest podprzestrzenią wektorową \(V\);
\item \(W\neq\varnothing\) i dla wszystkich \(u,v\in W\),
      \(a,b\in K\) zachodzi \(au+bv\in W\);
\item \(0_V\in W\), \(u+v\in W\) dla \(u,v\in W\)
      oraz \(au\in W\) dla \(a\in K\), \(u\in W\).
\end{enumerate}
\end{twierdzenie}
\begin{proof}
(i)\(\Rightarrow\)(iii): podprzestrzeń jest przestrzenią
wektorową; jej zero pokrywa się z zerem \(V\),
bo \(0_Kw=0_V\) dla każdego \(w\in W\), a \(W\) jest niepusta.
Domknięcie działań jest częścią definicji.
(iii)\(\Rightarrow\)(ii) wynika z kolejnego zastosowania
domknięcia mnożenia przez skalar i dodawania.

Zakładając (ii), wybieramy \(w\in W\).
Dla \(a=b=0\) i \(u=v=w\) dostajemy \(0_V\in W\).
Z (ii) otrzymujemy \(u+v\in W\) (dla \(a=b=1\))
oraz \(au\in W\) (dla \(b=0\)).
W szczególności \(-u=(-1)u\in W\).
Aksjomaty łączności, przemienności i rozdzielności
obowiązują w \(W\), ponieważ są prawdziwe w \(V\).
Otrzymujemy (i).
\end{proof}

\begin{przyklad}[Równania jednorodne i niejednorodne]
\label{prz:lin-rownania-jednorodne}
W \(\mathbb R^3\) rozważmy
\(W=\{(x,y,z):x-2y+z=0\}\).
Jeżeli \(u=(x_1,y_1,z_1)\) i \(v=(x_2,y_2,z_2)\)
spełniają równanie, to dla \(a,b\in\mathbb R\)
współrzędne \(au+bv\) spełniają
\(a(x_1-2y_1+z_1)+b(x_2-2y_2+z_2)=0\).
Zatem \(W\) jest podprzestrzenią.

Zbiór \(A=\{(x,y,z):x-2y+z=1\}\) jest niepusty,
ale nie zawiera \(0\), więc nie jest podprzestrzenią.
Jest przesuniętą płaszczyzną, czyli zbiorem afinicznym.
\end{przyklad}

\begin{figure}[!htbp]
\centering
\begin{tikzpicture}[>=Latex,scale=0.94,every node/.style={font=\small}]
\begin{scope}
 \draw[->,gray] (-1.55,0)--(1.7,0) node[right] {\(x\)};
 \draw[->,gray] (0,-1.0)--(0,1.30) node[above] {\(y\)};
 \draw[line width=1.8pt,manifoldblue] (-1.45,-0.72)--(1.45,0.72);
 \fill (0,0) circle (1.8pt);
 \node[manifoldblue,above left] at (-1.0,-0.50) {\(W\)};
 \node[below right] at (0,0) {\(0\)};
 \node[below] at (0,-1.25) {Podprzestrzeń};
\end{scope}
\begin{scope}[xshift=5.6cm]
 \draw[->,gray] (-1.55,0)--(1.7,0) node[right] {\(x\)};
 \draw[->,gray] (0,-1.0)--(0,1.30) node[above] {\(y\)};
 \draw[line width=1.8pt,manifoldred] (-1.45,-0.10)--(1.25,1.25);
 \fill (0,0) circle (1.8pt);
 \node[manifoldred,above left] at (-0.98,0.10) {\(A\)};
 \node[below right] at (0,0) {\(0\)};
 \node[below] at (0,-1.25) {Przesunięta prosta};
\end{scope}
\end{tikzpicture}
\caption{Ilustracja w \(\mathbb R^2\): prosta przechodząca
przez zero jest podprzestrzenią, a prosta przesunięta
nie jest podprzestrzenią. W \(\mathbb R^3\) analogicznie
rozróżniamy płaszczyzny liniowe i afiniczne.}
\label{fig:lin-podprzestrzen-afiniczna}
\end{figure}

\begin{twierdzenie}[Przecięcia podprzestrzeni]
\label{thm:lin-przeciecie}
Niech \((W_i)_{i\in I}\) będzie dowolną rodziną
podprzestrzeni \(V\). Wtedy \(\bigcap_{i\in I}W_i\)
jest podprzestrzenią \(V\). Dla \(I=\varnothing\)
przyjmujemy \(\bigcap_{i\in I}W_i=V\).
\end{twierdzenie}
\begin{proof}
Każde \(W_i\) zawiera \(0_V\), zatem przecięcie
jest niepuste. Jeśli \(u,v\) należą do wszystkich
\(W_i\), to dla \(a,b\in K\) wektor \(au+bv\)
należy do każdego \(W_i\), a więc również do przecięcia.
Stosujemy Twierdzenie~\ref{thm:lin-kryterium-podprzestrzeni}.
Przypadek pustej rodziny wynika z przyjętej konwencji.
\end{proof}

\begin{twierdzenie}[Kiedy unia podprzestrzeni jest podprzestrzenią?]
\label{thm:lin-unia}
Jeśli \(U,W\subseteq V\) są podprzestrzeniami,
to \(U\cup W\) jest podprzestrzenią dokładnie wtedy,
gdy \(U\subseteq W\) lub \(W\subseteq U\).
\end{twierdzenie}
\begin{proof}
Jeżeli jedna podprzestrzeń zawiera drugą, ich unia
jest większą podprzestrzenią.
Odwrotnie, załóżmy, że żadna nie zawiera drugiej.
Możemy wybrać \(u\in U\setminus W\) i \(w\in W\setminus U\).
Gdyby \(u+w\in U\), to \(w=(u+w)-u\in U\), sprzeczność.
Gdyby \(u+w\in W\), to \(u=(u+w)-w\in W\), również sprzeczność.
Zatem \(U\cup W\) nie jest zamknięta na dodawanie.
\end{proof}

\begin{twierdzenie}[Skończona unia podprzestrzeni nad ciałem nieskończonym]
\label{thm:lin-skonczona-unia-nieskonczone-cialo}
Niech \(K\) będzie ciałem nieskończonym,
\(V\) przestrzenią wektorową nad \(K\),
a \(U_1,\ldots,U_m\subsetneq V\) jej właściwymi
podprzestrzeniami, gdzie \(m\geq1\).
Wtedy \(U_1\cup\cdots\cup U_m\neq V\).
\end{twierdzenie}
\begin{proof}
Przypuśćmy przeciwnie, że \(V=U_1\cup\cdots\cup U_m\).
Spośród tych skończenie wielu podprzestrzeni
wybieramy podrodzinę pokrywającą \(V\)
o najmniejszej liczbie elementów i numerujemy
ją ponownie \(U_1,\ldots,U_r\).
Ponieważ każdy \(U_i\) jest właściwy, mamy \(r\geq2\).
Minimalność pokrycia oznacza, że dla każdego \(i\)
istnieje wektor należący do \(U_i\), ale do żadnego
z pozostałych \(U_j\): w przeciwnym razie można by
usunąć \(U_i\), nie tracąc pokrycia.
Wybierzmy zatem
\(x\in U_1\setminus\bigcup_{j=2}^r U_j\)
oraz \(y\in U_2\setminus\bigcup_{j\neq2} U_j\).

Rozważmy wektory \(x+ty\), gdzie
\(t\in K\setminus\{0\}\).
Żaden z nich nie należy do \(U_1\),
bo wtedy \(ty=(x+ty)-x\in U_1\),
a z \(t\neq0\) wynikałoby \(y\in U_1\).
Żaden nie należy też do \(U_2\):
gdyby \(x+ty\in U_2\), to
\(x=(x+ty)-ty\in U_2\), sprzeczność.

Dla każdego \(j\in\{3,\ldots,r\}\)
co najwyżej jeden skalar \(t\in K\)
może spełniać \(x+ty\in U_j\).
Istotnie, gdyby zachodziło to dla
\(t_1\neq t_2\), otrzymalibyśmy
\((t_1-t_2)y\in U_j\), więc \(y\in U_j\),
wbrew wyborowi \(y\).
Zatem spośród wszystkich niezerowych skalarów
tylko skończenie wiele może prowadzić
do wektora \(x+ty\) leżącego w którymkolwiek \(U_j\).
Ponieważ \(K\) jest nieskończone, istnieje
niezerowy \(t\) spoza tej skończonej rodziny.
Wtedy \(x+ty\notin\bigcup_{j=1}^rU_j=V\),
co jest niemożliwe.
\end{proof}

\begin{przyklad}[Dlaczego nieskończoność ciała jest istotna?]
\label{prz:lin-unia-trzech-f2}
W \(\mathbb F_2^2\) rozważmy trzy podprzestrzenie
\(U=\operatorname{span}\{(1,0)\}\),
\(W=\operatorname{span}\{(0,1)\}\) i
\(Z=\operatorname{span}\{(1,1)\}\).
Każda ma dwa elementy: wektor zerowy
oraz wskazany niezerowy generator.
Wektory niezerowe w \(\mathbb F_2^2\)
to dokładnie \((1,0),(0,1),(1,1)\).
Stąd \(U\cup W\cup Z=\mathbb F_2^2\),
choć żadna z trzech podprzestrzeni
nie zawiera pozostałych, a każda jest właściwa.
Twierdzenie~\ref{thm:lin-unia} o \emph{dwóch}
podprzestrzeniach obowiązuje nad każdym ciałem,
natomiast jego naiwne uogólnienie na
dowolną skończoną liczbę podprzestrzeni
jest fałszywe nad ciałami skończonymi.
\end{przyklad}

\begin{uwaga}[Nie mylić unii ze sumą podprzestrzeni]
\label{rem:lin-unia-suma}
Zbiory \(U\cup W\) i
\(U+W:=\{u+w:u\in U,\ w\in W\}\)
są na ogół różne. Ten drugi zbiór zawsze jest podprzestrzenią:
z \(u_1+w_1,u_2+w_2\in U+W\) dostajemy
\(a(u_1+w_1)+b(u_2+w_2)
=(au_1+bu_2)+(aw_1+bw_2)\in U+W\).
Własności sum i sum prostych rozwiniemy w Rozdziale~2;
na razie ważne jest odróżnienie sumy algebraicznej od unii.
\end{uwaga}

\section{Kombinacje liniowe i przestrzeń rozpięta}
\label{sec:lin-rozpietosc}

\begin{definicja}[Kombinacja liniowa i przestrzeń rozpięta]
\label{def:lin-span}
Niech \(S\subseteq V\). \emph{Kombinacją liniową}
wektorów z \(S\) nazywamy sumę
\(\sum_{j=1}^{m}a_js_j\), gdzie \(m\geq0\),
\(a_j\in K\) oraz \(s_j\in S\).
Suma pusta (przypadek \(m=0\)) jest równa \(0_V\).
\emph{Przestrzeń rozpięta} przez \(S\) jest zbiorem
wszystkich jego kombinacji liniowych:
\[
 \operatorname{span}_K(S)
   :=\left\{\sum_{j=1}^{m}a_js_j:
                    m\geq0,\ a_j\in K,\ s_j\in S\right\}.
\]
Piszemy \(\operatorname{span}(S)\), gdy ciało jest ustalone.
W szczególności \(\operatorname{span}(\varnothing)=\{0_V\}\).
Mówimy, że \(S\) \emph{rozpina} \(V\), jeżeli
\(\operatorname{span}(S)=V\).
\end{definicja}

Słowo \emph{skończona} jest tu niezbędne:
nawet jeżeli \(S\) jest nieskończony, pojedyncza
kombinacja liniowa wykorzystuje tylko skończenie wiele wektorów.
Nieskończony szereg wymagałby dodatkowego pojęcia zbieżności,
którego w ogólnej przestrzeni wektorowej nie ma.

\begin{twierdzenie}[Minimalna podprzestrzeń zawierająca zbiór]
\label{thm:lin-span-minimalny}
Dla dowolnego \(S\subseteq V\) zbiór
\(\operatorname{span}(S)\) jest najmniejszą względem
zawierania podprzestrzenią \(V\) zawierającą \(S\).
W szczególności
\[
 \operatorname{span}(S)
 =\bigcap\{W\subseteq V:W\text{ jest podprzestrzenią oraz }S\subseteq W\}.
\]
\end{twierdzenie}
\begin{proof}
Zbiór \(\operatorname{span}(S)\) zawiera \(0_V\), bo dopuszczamy
sumę pustą. Jeśli \(x=\sum_{j=1}^m a_js_j\)
i \(y=\sum_{k=1}^r b_kt_k\), to
\(cx+dy=\sum_{j=1}^m(ca_j)s_j+\sum_{k=1}^r(db_k)t_k\)
jest kombinacją liniową elementów \(S\).
Kryterium podprzestrzeni daje zatem, że
\(\operatorname{span}(S)\) jest podprzestrzenią.
Zawiera każdy \(s\in S\) jako kombinację \(1_Ks\).

Jeśli \(W\) jest dowolną podprzestrzenią zawierającą \(S\),
to z domknięcia na dodawanie i mnożenie przez skalary
wynika, że \(W\) zawiera każdą skończoną kombinację liniową
elementów \(S\). Zatem \(\operatorname{span}(S)\subseteq W\).
Powyższa rodzina podprzestrzeni jest niepusta, bo zawiera \(V\),
więc to samo rozumowanie daje żądaną równość.
\end{proof}

\begin{stwierdzenie}[Podstawowe rachunki na przestrzeniach rozpiętych]
\label{prop:lin-span-reguly}
Dla \(A,B\subseteq V\) zachodzą:
\(\operatorname{span}(A)\subseteq\operatorname{span}(B)\),
jeśli \(A\subseteq B\);
\(\operatorname{span}(\operatorname{span}(A))
=\operatorname{span}(A)\);
oraz
\(\operatorname{span}(A\cup B)
=\operatorname{span}(A)+\operatorname{span}(B)\).
\end{stwierdzenie}
\begin{proof}
Każda kombinacja elementów \(A\) jest kombinacją
elementów \(B\), co dowodzi pierwszej własności.
Druga wynika z faktu, że \(\operatorname{span}(A)\)
jest podprzestrzenią. Dla trzeciej rozdzielamy
każdą skończoną kombinację elementów \(A\cup B\)
na sumę kombinacji elementów \(A\) i elementów \(B\).
Daje to zawieranie z lewej do prawej. Z drugiej strony,
suma dwóch takich kombinacji jest kombinacją elementów
\(A\cup B\), co daje zawieranie przeciwne.
\end{proof}

\begin{przyklad}[Opis płaszczyzny przez generatory]
\label{prz:lin-span-r3}
Weźmy \(u=(1,0,1)\) i \(w=(0,1,1)\) w \(\mathbb R^3\).
Każdy wektor ich przestrzeni rozpiętej ma postać
\(au+bw=(a,b,a+b)\), a więc spełnia \(z=x+y\).
Odwrotnie, jeżeli \(z=x+y\), to
\((x,y,z)=xu+yw\). Otrzymaliśmy dokładny opis
\(\operatorname{span}\{u,w\}
=\{(x,y,z)\in\mathbb R^3:z=x+y\}\).
Wektor \((1,1,0)\) nie należy do tej płaszczyzny,
ponieważ \(0\neq1+1\).
\end{przyklad}

\section{Niezależność i zależność liniowa}
\label{sec:lin-niezaleznosc}

Rozpinanie odpowiada na pytanie, \emph{czy istnieje}
przedstawienie danego wektora przez wybrane generatory.
Niezależność odpowiada na inne pytanie: \emph{czy współczynniki
takiego przedstawienia są jednoznaczne}.

\begin{definicja}[Niezależność liniowa]
\label{def:lin-niezaleznosc}
Rodzinę wektorów \((v_i)_{i\in I}\) nazywamy
\emph{liniowo niezależną}, jeżeli dla dowolnych parami
różnych \(i_1,\ldots,i_m\in I\) i skalarów
\(a_1,\ldots,a_m\in K\) równość
\(a_1v_{i_1}+\cdots+a_mv_{i_m}=0_V\)
pociąga \(a_1=\cdots=a_m=0_K\).
Jeśli ten warunek nie zachodzi, rodzina jest
\emph{liniowo zależna}.
Pusta rodzina jest liniowo niezależna.
\end{definicja}

Niezależność dotyczy \emph{rodziny indeksowanej}.
W szczególności ta sama wartość wektora wpisana dwa razy
pod różnymi indeksami tworzy rodzinę zależną, jeśli \(1_K\neq0_K\):
\(v-v=0\). Dla zbioru \(S\subseteq V\) stosujemy
tę definicję do rodziny jego różnych elementów
z indeksami \(s\in S\).

\begin{lemat}[Równoważne kryteria niezależności]
\label{lem:lin-niezaleznosc-kryteria}
Niech \((v_i)_{i\in I}\) będzie rodziną wektorów w \(V\).
Następujące warunki są równoważne:
\begin{enumerate}[label=(\roman*)]
\item rodzina jest liniowo niezależna;
\item każdy \(x\in\operatorname{span}\{v_i:i\in I\}\)
      ma dokładnie jeden zapis w postaci
      \(\sum_{i\in I}a_iv_i\), gdzie rodzina
      współczynników \((a_i)\) ma skończony nośnik;
\item dla każdego \(j\in I\) wektor \(v_j\)
      nie należy do \(\operatorname{span}\{v_i:i\in I,\ i\neq j\}\).
\end{enumerate}
\end{lemat}
\begin{proof}
(i)\(\Rightarrow\)(ii): istnienie zapisu wynika z definicji
przestrzeni rozpiętej. Jeśli dwa zapisy przedstawiają
ten sam wektor, odejmujemy je i zbieramy współczynniki
przy jednakowych indeksach. Ponieważ suma ich nośników
jest skończona, otrzymujemy skończoną kombinację równą zeru;
niezależność wymusza równość wszystkich odpowiadających
współczynników. Zatem zapisy są identyczne.
(ii)\(\Rightarrow\)(i): wektor zerowy ma zapis z wszystkimi
współczynnikami zerowymi, więc każdy inny zapis o
skończonym nośniku byłby sprzeczny z jednoznacznością.

(i)\(\Rightarrow\)(iii): równość
\(v_j=\sum_{i\neq j}a_iv_i\) dawałaby
nietrywialną zależność ze współczynnikiem \(1_K\)
przy \(v_j\).
(iii)\(\Rightarrow\)(i): jeśli istnieje nietrywialna
skończona relacja \(\sum a_iv_i=0\), wybieramy
\(j\) z \(a_j\neq0_K\). Po pomnożeniu przez
\(a_j^{-1}\) otrzymujemy
\(v_j=-a_j^{-1}\sum_{i\neq j}a_iv_i\),
sprzecznie z (iii).
\end{proof}

\begin{wniosek}[Podrodziny i dołączanie wektora]
\label{cor:lin-podrodziny}
Każda podrodzina rodziny liniowo niezależnej jest
liniowo niezależna. Jeżeli \(S\subseteq V\)
jest liniowo niezależnym zbiorem, a \(v\in V\setminus S\),
to \(S\cup\{v\}\) jest liniowo niezależny
dokładnie wtedy, gdy \(v\notin\operatorname{span}(S)\).
\end{wniosek}
\begin{proof}
Relacja liniowa między elementami podrodziny
jest także relacją między elementami większej rodziny.
Jeśli \(v\in\operatorname{span}(S)\), to istnieją
\(s_1,\ldots,s_m\in S\) i \(a_1,\ldots,a_m\in K\)
takie, że \(v=\sum_{j=1}^m a_js_j\).
Ponieważ \(v\notin S\), relacja
\(v-\sum_{j=1}^m a_js_j=0\)
jest nietrywialna: współczynnik przy \(v\)
wynosi \(1_K\). Zatem \(S\cup\{v\}\) jest zależny.

Odwrotnie, niech \(v\notin\operatorname{span}(S)\).
W dowolnej relacji
\(av+\sum_{j=1}^m a_js_j=0\), \(s_j\in S\),
zbieramy współczynniki przy powtarzających się
elementach zbioru \(S\). Jeśli \(a\neq0_K\),
to \(v=-a^{-1}\sum_{j=1}^m a_js_j\)
należałby do \(\operatorname{span}(S)\), sprzeczność.
Zatem \(a=0_K\), a niezależność \(S\)
wymusza zerowanie pozostałych współczynników.
\end{proof}

\begin{uwaga}[Zbiory a rodziny]
\label{rem:lin-zbior-rodzina}
Gdy \(v\in S\), unia zbiorów \(S\cup\{v\}=S\)
pozostaje niezależna, choć \(v\in\operatorname{span}(S)\).
Dlatego w kryterium dla \emph{zbiorów}
założyliśmy \(v\notin S\).
Dla \emph{rodziny indeksowanej} dołączenie wektora
pod nowym indeksem oznacza nowe wystąpienie;
taka rozszerzona rodzina jest niezależna
dokładnie wtedy, gdy \(v\notin\operatorname{span}(S)\).
Jeżeli wartość \(v\) występowała już pod innym
indeksem, rozszerzona rodzina jest zależna
z powodu relacji \(v-v=0\) między tymi indeksami.
\end{uwaga}

\begin{przyklad}[Trzy wektory na płaszczyźnie]
\label{prz:lin-niezaleznosc-r2}
W \(\mathbb R^2\) wektory \(e_1=(1,0)\), \(e_2=(0,1)\)
są niezależne: z \(ae_1+be_2=(a,b)=0\)
wynika \(a=b=0\).
Po dołączeniu \(w=(1,1)\) powstaje rodzina zależna,
bo \(e_1+e_2-w=0\). Żaden pojedynczy wektor
spośród \(e_1,e_2,w\) nie jest zerowy; zależność
wynika wyłącznie z relacji między nimi.
\end{przyklad}

\begin{figure}[!htbp]
\centering
\begin{tikzpicture}[>=Latex,scale=0.93,every node/.style={font=\small}]
 \begin{scope}
  \draw[->,gray!70] (-0.45,0)--(2.35,0) node[right] {\(x\)};
  \draw[->,gray!70] (0,-0.3)--(0,2.0) node[above] {\(y\)};
  \draw[->,very thick,manifoldblue] (0,0)--(1.9,0)
       node[below right] {\(e_1\)};
  \draw[->,very thick,manifoldgreen] (0,0)--(0,1.6)
       node[above left] {\(e_2\)};
  \node at (0.95,-0.76) {Dwa wektory niezależne};
 \end{scope}
 \begin{scope}[xshift=5.9cm]
  \draw[->,gray!70] (-0.45,0)--(2.35,0) node[right] {\(x\)};
  \draw[->,gray!70] (0,-0.3)--(0,2.0) node[above] {\(y\)};
  \draw[densely dashed,gray!70] (1.9,0)--(1.9,1.6)--(0,1.6);
  \draw[->,very thick,manifoldblue] (0,0)--(1.9,0)
       node[below right] {\(e_1\)};
  \draw[->,very thick,manifoldgreen] (0,0)--(0,1.6)
       node[above left] {\(e_2\)};
  \draw[->,very thick,manifoldred] (0,0)--(1.9,1.6)
       node[above right] {\(w=e_1+e_2\)};
  \node at (1.0,-0.76) {Trzy wektory zależne};
 \end{scope}
\end{tikzpicture}
\caption{Niezależność nie oznacza jedynie niezerowości wektorów.
Dwa wektory \(e_1,e_2\) są niezależne, lecz po dołączeniu
\(w=e_1+e_2\) otrzymujemy rodzinę zależną.
Linie przerywane ilustrują dodawanie wektorów.}
\label{fig:lin-niezaleznosc-r2}
\end{figure}

\begin{przyklad}[Charakterystyka dwa]
\label{prz:lin-f2}
W \(\mathbb F_2^2\) występują trzy niezerowe wektory:
\(u=(1,0)\), \(v=(0,1)\), \(w=(1,1)\).
Każde dwa różne wektory niezerowe są niezależne.
Istotnie, jeśli dla takich wektorów \(a,b\)
mamy \(\alpha a+\beta b=0\) z
\(\alpha,\beta\in\mathbb F_2\), to tylko
współczynniki \(0\) i \(1\) są możliwe.
Przypadki \((\alpha,\beta)=(1,0)\)
oraz \((0,1)\) przeczą niezerowości
wektorów. Przypadek \((1,1)\)
dawałby \(a+b=0\), a więc \(a=b\)
w charakterystyce \(2\), również sprzeczność.
Stąd \(\alpha=\beta=0\).
Jednak wszystkie trzy wektory są zależne, ponieważ
\(u+v+w=(0,0)\) nad \(\mathbb F_2\).
W \(\mathbb F_2^2\) nie można więc opisywać
niezależności za pomocą kątów ani długości.
\end{przyklad}

\section{Bazy i wymiar w przypadku skończonym}
\label{sec:lin-bazy-skonczone}

\begin{definicja}[Baza i wymiar]
\label{def:lin-baza}
\emph{Bazą} przestrzeni \(V\) nazywamy rodzinę
wektorów, która jest liniowo niezależna i rozpina \(V\).
Przestrzeń nazywamy \emph{skończenie wymiarową},
jeżeli posiada skończoną bazę.
W szczególności pusta rodzina jest bazą przestrzeni \(\{0\}\).
Jeżeli \(V\) ma skończoną bazę składającą się z \(n\)
wektorów, to liczbę \(n\) nazywamy jej \emph{wymiarem}
i oznaczamy \(\dim_KV=n\).
Wykażemy niżej, że liczba ta nie zależy od wybranej bazy.
\end{definicja}

Baza łączy dwa warunki, które oddzielnie nie wystarczają:
generatory mogą zawierać powtórzenia i nadmiarowe wektory,
natomiast rodzina niezależna może nie rozpinać całej przestrzeni.

\begin{lemat}[Usuwanie zbędnych generatorów]
\label{lem:lin-usuwanie-generatorow}
Jeżeli skończona rodzina \((s_1,\ldots,s_m)\)
rozpina \(V\), to można usunąć z niej pewne
wektory, otrzymując bazę \(V\).
\end{lemat}
\begin{proof}
Jeśli rodzina jest niezależna, jest już bazą.
W przeciwnym razie istnieją skalary
\(a_1,\ldots,a_m\), nie wszystkie zerowe, takie że
\(\sum_{j=1}^m a_js_j=0\).
Wybieramy \(k\) z \(a_k\neq0\). Wówczas
\(s_k=-a_k^{-1}\sum_{j\neq k}a_js_j\).
Każde wystąpienie \(s_k\) w kombinacji generatorów
można zastąpić tą sumą, więc rodzina bez \(s_k\)
nadal rozpina \(V\). Powtarzamy procedurę.
Liczba wektorów za każdym razem maleje, toteż
po co najwyżej \(m\) krokach otrzymamy rodzinę
niezależną. Przypadek \(V=\{0\}\) kończy się
rodziną pustą, która nadal rozpina \(V\).
\end{proof}

\begin{twierdzenie}[Lemat o wymianie Steinitza]
\label{thm:lin-steinitz}
Niech \(V\) będzie przestrzenią wektorową nad \(K\).
Załóżmy, że \(V\) rozpina skończona rodzina
\((s_1,\ldots,s_m)\) i że
\((v_1,\ldots,v_r)\) jest rodziną liniowo niezależną.
Wtedy \(r\leq m\). Co więcej, można wybrać
\(r\) pozycji w rodzinie \((s_1,\ldots,s_m)\)
i zastąpić znajdujące się w nich wektory kolejno
przez \(v_1,\ldots,v_r\), tak by powstała
\(m\)-elementowa rodzina rozpinająca \(V\).
Nie zakładamy, że początkowa rodzina generatorów
jest niezależna.
\end{twierdzenie}
\begin{proof}
Najpierw, dla \(0\leq k\leq\min\{m,r\}\),
konstruujemy indukcyjnie rodzinę rozpinającą, złożoną
z \(v_1,\ldots,v_k\) oraz \(m-k\)
wektorów pozostawionych z wyjściowej rodziny.
Przy \(k=0\) jest to założenie.

Załóżmy, że zbudowaliśmy taką rodzinę
dla \(k<\min\{m,r\}\).
Ponieważ rozpina ona \(V\), wektor \(v_{k+1}\)
ma przedstawienie
\(v_{k+1}=\sum_{i=1}^{k}a_iv_i+
                 \sum_{j\in J_k}b_js_j\),
gdzie \(J_k\) jest zbiorem \(m-k\)
pozostałych pozycji generatorów.
Co najmniej jeden z \(b_j\) jest niezerowy.
Gdyby wszystkie były zerowe, wektor \(v_{k+1}\)
należałby do przestrzeni rozpiętej przez
\(v_1,\ldots,v_k\), co przeczy niezależności
rodziny \(v_1,\ldots,v_r\).

Wybieramy więc \(j_0\in J_k\) z
\(b_{j_0}\neq0\) i rozwiązujemy tę równość
względem \(s_{j_0}\):
\[
 s_{j_0}=b_{j_0}^{-1}v_{k+1}
   -\sum_{i=1}^{k}b_{j_0}^{-1}a_iv_i
   -\sum_{j\in J_k\setminus\{j_0\}}
                       b_{j_0}^{-1}b_js_j .
\]
Zatem nowa rodzina, w której \(s_{j_0}\)
zastąpiono \(v_{k+1}\), rozpina wszystkie
wektory poprzedniej rodziny, a więc całe \(V\).
Ma \(k+1\) wektorów z rodziny \(v_i\)
i \(m-k-1\) pozostałych generatorów.
Zachowujemy tezę indukcyjną.

Jeżeli \(r>m\), to po \(m\) krokach
otrzymujemy rodzinę \((v_1,\ldots,v_m)\)
rozpinającą \(V\), ponieważ wykorzystaliśmy
wszystkie \(m\) pozycji generatorów.
Wówczas \(v_{m+1}\) również należałby
do \(\operatorname{span}\{v_1,\ldots,v_m\}\),
wbrew niezależności rodziny \((v_1,\ldots,v_r)\).
Zatem \(r\leq m\). Indukcję można więc
przeprowadzić do kroku \(r\), otrzymując
zapowiedzianą rodzinę rozpinającą.
\end{proof}

\begin{wniosek}[Jednakowa liczba elementów baz skończonych]
\label{cor:lin-bazy-rowna-liczba}
Jeżeli \(B=(b_1,\ldots,b_n)\) i
\(C=(c_1,\ldots,c_m)\) są skończonymi bazami
tej samej przestrzeni \(V\), to \(n=m\).
\end{wniosek}
\begin{proof}
Rodzina \(B\) jest niezależna, a \(C\) rozpina \(V\),
więc z Twierdzenia~\ref{thm:lin-steinitz}
wynika \(n\leq m\).
Zamieniając role \(B\) i \(C\), dostajemy \(m\leq n\).
Stąd \(n=m\), co uzasadnia jednoznaczność
wymiaru z Definicji~\ref{def:lin-baza}.
\end{proof}

\begin{twierdzenie}[Uzupełnianie rodziny niezależnej do bazy]
\label{thm:lin-uzupelnianie-skonczone}
Niech \(V\) będzie rozpięta przez skończoną rodzinę
\((s_1,\ldots,s_m)\), a
\((v_1,\ldots,v_r)\) niech będzie rodziną
liniowo niezależną w \(V\).
Wówczas pewne wektory spośród \(s_1,\ldots,s_m\)
można dołączyć do \((v_1,\ldots,v_r)\)
tak, aby otrzymać bazę \(V\).
\end{twierdzenie}
\begin{proof}
Rozważamy skończoną rodzinę
\((v_1,\ldots,v_r,s_1,\ldots,s_m)\).
Przechodzimy kolejno przez elementy
\(s_1,\ldots,s_m\). Wektor \(s_j\) dołączamy
do dotychczas wybranej rodziny dokładnie wtedy,
gdy nie leży w jej przestrzeni rozpiętej.
Ponieważ dołączamy \(s_j\) tylko wtedy, gdy
nie należy do przestrzeni rozpiętej przez już
wybrane wektory, nie jest on w szczególności
powtórzeniem żadnego z nich.
Z Wniosku~\ref{cor:lin-podrodziny}
wynika więc, że po każdym kroku
zachowujemy liniową niezależność.
Jeśli wektor pominięto, był już kombinacją
wektorów dotychczas wybranych, więc także
końcowych. Jeśli dołączono, również należy
do ich przestrzeni rozpiętej.
Po ostatnim kroku wszystkie \(s_j\)
leżą więc w przestrzeni rozpiętej przez wybraną
rodzinę. Skoro \(s_j\) rozpinają \(V\), wybrana
rodzina też rozpina \(V\), a zatem jest bazą.
\end{proof}

\begin{wniosek}[Kryteria bycia bazą]
\label{cor:lin-baza-kryteria}
Niech \(V\) ma wymiar \(n<\infty\).
Wtedy:
\begin{enumerate}[label=(\roman*)]
\item każda rodzina niezależna ma co najwyżej \(n\) wektorów;
\item każda skończona rodzina rozpinająca ma co najmniej \(n\) wektorów;
\item rodzina złożona z dokładnie \(n\) wektorów jest bazą,
      jeżeli jest niezależna \emph{albo} rozpina \(V\).
\end{enumerate}
\end{wniosek}
\begin{proof}
W punkcie (i) stosujemy Twierdzenie~\ref{thm:lin-steinitz}
do dowolnej niezależnej rodziny i ustalonej bazy
\(n\)-elementowej. Gdyby niezależna rodzina
miała więcej niż \(n\) elementów, jej podrodzina
z \(n+1\) elementami przeczyłaby temu twierdzeniu.
Punkt (ii) otrzymujemy, stosując to samo twierdzenie
do ustalonej bazy niezależnej i rodziny rozpinającej.

Jeżeli rodzina \(n\)-elementowa jest niezależna,
uzupełniamy ją do bazy z
Twierdzenia~\ref{thm:lin-uzupelnianie-skonczone}.
Nowa baza musi mieć \(n\) elementów,
więc nic nie można było dołączyć;
rodzina wyjściowa już rozpina \(V\).
Jeżeli rodzina \(n\)-elementowa rozpina \(V\),
usuwamy z niej zależne elementy
zgodnie z Lematem~\ref{lem:lin-usuwanie-generatorow}.
Powstała baza ma \(n\) elementów,
więc żaden element nie został usunięty.
\end{proof}

\begin{twierdzenie}[Wymiar podprzestrzeni]
\label{thm:lin-wymiar-podprzestrzeni}
Niech \(V\) będzie przestrzenią skończenie wymiarową
nad \(K\), a \(W\subseteq V\) jej podprzestrzenią.
Wtedy \(W\) jest skończenie wymiarowa oraz
\(\dim_KW\leq\dim_KV\).
Ponadto równość wymiarów zachodzi dokładnie wtedy,
gdy \(W=V\).
\end{twierdzenie}
\begin{proof}
Oznaczmy \(n=\dim_KV\). Jeśli \(W=\{0\}\),
to pusta rodzina jest bazą \(W\), a więc
\(\dim_KW=0\leq n\).
Załóżmy, że \(W\neq\{0\}\).
Wybierzmy \(w_1\in W\setminus\{0\}\).
Rodzina \((w_1)\) jest niezależna.
Jeśli \(\operatorname{span}\{w_1\}\neq W\),
wybieramy \(w_2\in W\setminus\operatorname{span}\{w_1\}\).
Z Wniosku~\ref{cor:lin-podrodziny} wynika,
że \((w_1,w_2)\) jest niezależna.
Postępujemy tak dalej: jeżeli wybrane wektory
nie rozpinają \(W\), dołączamy jeden z \(W\)
spoza ich przestrzeni rozpiętej.
Każdy krok powiększa niezależną rodzinę.
Ponieważ taka rodzina w \(V\) ma co najwyżej
\(n\) elementów (Wniosek~\ref{cor:lin-baza-kryteria}),
procedura musi zakończyć się po skończenie wielu krokach.
Otrzymujemy bazę \(W\) liczącą \(r\leq n\) wektorów.

Jeżeli \(W=V\), równość wymiarów jest oczywista.
Odwrotnie, jeśli \(r=n\), otrzymana baza
\(W\) jest niezależną rodziną \(n\) wektorów w \(V\).
Z Wniosku~\ref{cor:lin-baza-kryteria}(iii)
rozpina ona także \(V\). Wszystkie jej
wektory należą do \(W\), więc \(V\subseteq W\),
a stąd \(V=W\).
\end{proof}

\begin{przyklad}[Bazy kanoniczne przestrzeni współrzędnych i macierzy]
\label{prz:lin-bazy-macierzy}
Niech \(n\geq1\). Dla \(j\in\{1,\ldots,n\}\)
oznaczmy przez \(e_j\in K^n\) wektor,
którego \(j\)-ta współrzędna wynosi \(1_K\),
a pozostałe \(0_K\). Każdy \(x=(x_1,\ldots,x_n)\)
ma zapis \(x=\sum_{j=1}^n x_je_j\);
jeżeli \(\sum_{j=1}^n a_je_j=0\),
to porównanie \(j\)-tej współrzędnej daje \(a_j=0\).
Zatem \((e_1,\ldots,e_n)\) jest bazą \(K^n\)
i \(\dim_K K^n=n\). Dla \(n=0\) odpowiednie
stwierdzenie pozostaje prawdziwe z bazą pustą.

Dla \(m,n\geq1\) niech
\(\operatorname{Mat}_{m\times n}(K)\) oznacza
przestrzeń wszystkich macierzy o \(m\) wierszach
i \(n\) kolumnach, z dodawaniem i mnożeniem
przez skalar wykonywanymi po elementach.
Niech \(E_{ij}\) będzie macierzą, której
wyraz na pozycji \((i,j)\) wynosi \(1_K\),
a pozostałe \(0_K\).
Każda macierz \(A=(a_{ij})\) ma zapis
\(A=\sum_{i=1}^m\sum_{j=1}^n a_{ij}E_{ij}\).
Jeśli ta suma jest macierzą zerową,
porównanie jej wyrazu \((i,j)\) daje \(a_{ij}=0\).
Macierze \(E_{ij}\) tworzą zatem bazę
i \(\dim_K\operatorname{Mat}_{m\times n}(K)=mn\).
\end{przyklad}

\begin{uwaga}[Definicja wymiaru a liczność rodziny]
\label{rem:lin-dim-nie-baza}
Dla przestrzeni \(\{0\}\) wymiar wynosi \(0\),
ponieważ pusta rodzina jest bazą. Nie istnieje
jednoelementowa baza przestrzeni zerowej:
jej jedyny wektor \(0\) jest liniowo zależny.
Nie wolno również wnioskować o wymiarze
wyłącznie z liczby podanych generatorów
bez zbadania ich niezależności.
\end{uwaga}

\section{Współrzędne względem bazy}
\label{sec:lin-wspolrzedne}

\begin{twierdzenie}[Jednoznaczność współrzędnych]
\label{thm:lin-wspolrzedne}
Niech \(V\) będzie \(n\)-wymiarową przestrzenią
nad \(K\), a \(B=(b_1,\ldots,b_n)\) jej bazą.
Dla każdego \(v\in V\) istnieje dokładnie jedna
kolumna \((a_1,\ldots,a_n)\in K^n\) taka, że
\(v=\sum_{i=1}^n a_ib_i\).
Odwzorowanie
\[
 [\,\cdot\,]_B\colon V\longrightarrow K^n,\qquad
 v\longmapsto(a_1,\ldots,a_n)
\]
jest bijekcją oraz zachowuje dodawanie
i mnożenie przez skalary.
\end{twierdzenie}
\begin{proof}
Istnienie zapisu wynika z faktu, że baza rozpina \(V\),
a jednoznaczność z
Lematu~\ref{lem:lin-niezaleznosc-kryteria}.
Jeśli \(v=\sum a_ib_i\) i \(w=\sum c_ib_i\),
to z rozdzielności działań wynika
\(v+w=\sum(a_i+c_i)b_i\) i
\(\lambda v=\sum(\lambda a_i)b_i\).
Jednoznaczność współczynników daje
\([v+w]_B=[v]_B+[w]_B\) oraz
\([\lambda v]_B=\lambda[v]_B\).
Dla dowolnej kolumny \((a_1,\ldots,a_n)\in K^n\)
wektor \(\sum a_ib_i\) ma dokładnie tę kolumnę
współrzędnych, więc odwzorowanie jest surjektywne.
Jeśli \([v]_B=[w]_B\), obie sumy są identyczne,
stąd \(v=w\), co dowodzi iniektywności.
\end{proof}

\begin{przyklad}[Współrzędne w niestandardowej bazie]
\label{prz:lin-wspolrzedne-niekanoniczne}
W \(\mathbb R^3\) rozważmy wektory
\(b_1=(1,1,0)\), \(b_2=(0,1,1)\),
\(b_3=(0,0,1)\).
Szukając współczynników \(a,b,c\) z
\((x,y,z)=ab_1+bb_2+cb_3\),
porównujemy współrzędne i otrzymujemy
\(x=a,\ y=a+b,\ z=b+c\).
Stąd jedyne możliwe rozwiązanie to
\(a=x\), \(b=y-x\), \(c=z-y+x\).
Jest ono poprawne dla każdego \((x,y,z)\in\mathbb R^3\),
a zatem rodzina \((b_1,b_2,b_3)\) jest bazą.
Na przykład
\([(2,3,5)]_B=(2,1,4)\), ponieważ
\(2b_1+b_2+4b_3=(2,3,5)\).
\end{przyklad}

\begin{przyklad}[Baza wielomianów]
\label{prz:lin-baza-wielomianow}
W \(K[X]_{\leq d}\) wektory \(1,X,\ldots,X^d\)
rozpinają przestrzeń z definicji wielomianu.
Równość \(\sum_{j=0}^d a_jX^j=0\)
jest równością wielomianów, a więc wszystkie
współczynniki \(a_j\) są zerowe.
Wektory są niezależne i stanowią bazę,
zatem \(\dim_KK[X]_{\leq d}=d+1\).
Dla \(p(X)=2-3X+X^3\) nad \(\mathbb R\)
w bazie \((1,X,X^2,X^3)\) mamy
\([p]_B=(2,-3,0,1)\).
\end{przyklad}

\begin{przyklad}[Współrzędne zależą od ciała]
\label{prz:lin-wspolrzedne-ciala}
Dla \(z=2+3i\in\mathbb C\)
w bazie \(B=(1,i)\) przestrzeni nad
\(\mathbb R\) mamy \([z]_B=(2,3)\).
W bazie \((1)\) przestrzeni \(\mathbb C\)
nad \(\mathbb C\) współrzędną jest jednoelementowa
kolumna \((2+3i)\).
Są to opisy tego samego elementu w dwóch różnych
przestrzeniach wektorowych (nad różnymi ciałami).
\end{przyklad}

\section{Bazy algebraiczne w dowolnym wymiarze}
\label{sec:lin-bazy-nieskonczone}

Definicje przestrzeni rozpiętej, niezależności i bazy
nie wymagają skończoności wymiaru.
W nieskończonej bazie \((b_i)_{i\in I}\)
każdy wektor ma \emph{skończony} zapis
\(\sum_{i\in F}a_ib_i\) z \(F\subseteq I\) skończonym.
Taką bazę nazywamy \emph{bazą Hamela}.
Nie należy jej mylić z bazą Schauderowską,
w której dopuszcza się granice nieskończonych sum;
to drugie pojęcie wymaga topologii i należy do analizy funkcjonalnej.

\begin{przyklad}[Baza kanoniczna ciągów o skończonym nośniku]
\label{prz:lin-baza-skonczony-nosnik}
Dla \(n\in\mathbb N_0\) niech \(e_n\) oznacza ciąg,
którego \(n\)-ty wyraz jest równy \(1_K\),
a pozostałe są zerowe.
Rodzina \((e_n)_{n\in\mathbb N_0}\)
jest niezależna: w równości
\(\sum_{n\in F}a_ne_n=0\) odczytanie współrzędnej
o indeksie \(n\in F\) daje \(a_n=0\).
Rozpina \(K^{(\mathbb N_0)}\), bo każdy ciąg
o skończonym nośniku \(F\) ma postać
\(\sum_{n\in F}x_ne_n\).

Ta sama rodzina \emph{nie} rozpina
\(K^{\mathbb N_0}\): każda skończona kombinacja
wektorów \(e_n\) ma skończony nośnik, natomiast
ciąg \((1,1,\ldots)\) ma nośnik nieskończony.
Formalny napis \(\sum_{n=0}^{\infty}e_n\)
nie jest kombinacją liniową w sensie
Definicji~\ref{def:lin-span}.
\end{przyklad}

\begin{przyklad}[Baza wszystkich wielomianów]
\label{prz:lin-baza-kx}
Rodzina \((1,X,X^2,\ldots)\) jest bazą Hamela \(K[X]\).
Każdy wielomian zawiera skończenie wiele potęg \(X\)
i jego zapis przez te potęgi jest jednoznaczny.
Przestrzeń \(K[X]\) nie ma skończonej bazy:
jeśli rozpinają ją \(p_1,\ldots,p_m\),
to przy istnieniu choć jednego niezerowego
generatora każda niezerowa ich kombinacja
ma stopień nie większy niż
\(d=\max\{\deg p_j:p_j\neq0\}\).
Zbiór, którego maksimum bierzemy, jest
wtedy niepusty i skończony.
Wtedy \(X^{d+1}\) nie należy do ich przestrzeni
rozpiętej. Jeśli wszystkie \(p_j\) są zerowe
(lub \(m=0\)), nie rozpinają nawet wielomianu \(1\).
\end{przyklad}

\begin{lemat}[Łańcuch rodzin niezależnych]
\label{lem:lin-lancuchy}
Niech \(\mathcal C\) będzie rodziną liniowo
uporządkowaną przez zawieranie,
której elementy są liniowo niezależnymi
podzbiorami \(V\).
Wtedy \(\bigcup_{S\in\mathcal C}S\)
jest liniowo niezależny.
\end{lemat}
\begin{proof}
W każdej ewentualnej relacji liniowej występuje
tylko skończenie wiele wektorów \(v_1,\ldots,v_m\).
Każdy należy do pewnego zbioru \(S_j\in\mathcal C\).
Ponieważ \(\mathcal C\) jest łańcuchem, spośród
skończonej rodziny \(S_1,\ldots,S_m\)
można wybrać zbiór największy względem zawierania.
Zawiera on wszystkie \(v_j\), a jest niezależny.
Stąd współczynniki relacji są zerowe.
Dla pustego łańcucha unia jest zbiorem pustym,
również niezależnym.
\end{proof}

\begin{twierdzenie}[Istnienie bazy i rozszerzanie rodziny niezależnej]
\label{thm:lin-istnienie-bazy-zorn}
Przyjmijmy aksjomat wyboru (równoważnie: lemat Zorna).
Dla każdej przestrzeni wektorowej \(V\) nad \(K\)
i każdego liniowo niezależnego zbioru
\(S_0\subseteq V\) istnieje baza Hamela \(B\)
taka, że \(S_0\subseteq B\).
W szczególności każda przestrzeń wektorowa ma bazę.
\end{twierdzenie}
\begin{proof}
Niech \(\mathcal P\) będzie rodziną wszystkich
liniowo niezależnych podzbiorów \(V\)
zawierających \(S_0\), uporządkowaną
przez zawieranie. Rodzina jest niepusta,
bo \(S_0\in\mathcal P\).
Niech \(\mathcal C\subseteq\mathcal P\)
będzie dowolnym łańcuchem.
Jeżeli \(\mathcal C=\varnothing\),
ograniczeniem górnym w \(\mathcal P\)
jest \(S_0\).
Jeżeli \(\mathcal C\neq\varnothing\),
jego unia \(\bigcup_{S\in\mathcal C}S\)
zawiera \(S_0\), ponieważ każdy
\(S\in\mathcal C\) zawiera \(S_0\).
Unia jest niezależna na mocy
Lematu~\ref{lem:lin-lancuchy},
więc należy do \(\mathcal P\)
i stanowi ograniczenie górne łańcucha.

Lemat Zorna zapewnia istnienie maksymalnego
elementu \(B\in\mathcal P\).
Przypuśćmy, że \(\operatorname{span}(B)\neq V\)
i wybierzmy \(v\in V\setminus\operatorname{span}(B)\).
Dowód Wniosku~\ref{cor:lin-podrodziny}
pokazuje, że \(B\cup\{v\}\) jest niezależny.
Zawiera ściśle \(B\) i nadal należy do
\(\mathcal P\), w sprzeczności z maksymalnością.
Zatem \(\operatorname{span}(B)=V\),
co oznacza, że \(B\) jest bazą.
\end{proof}

\begin{uwaga}[Zakres użycia aksjomatu wyboru]
\label{rem:lin-aksjomat-wyboru}
Dowód istnienia bazy dla przestrzeni skończenie
generowanych nie wymaga lematu Zorna:
wystarcza procedura z
Lematu~\ref{lem:lin-usuwanie-generatorow}.
Dla ogólnej przestrzeni nieskończenie wymiarowej
w powyższym dowodzie zastosowaliśmy aksjomat wyboru.
Bez dodatkowych założeń teoriomnogościowych nie
można w ZF dowieść ogólnego twierdzenia,
że każda przestrzeń wektorowa ma bazę.
Jednoznaczność mocy dwóch baz nieskończonych
będzie tematem następnego rozdziału,
gdzie omówimy wymiar jako liczbę kardynalną.
\end{uwaga}

\begin{twierdzenie}[Charakteryzacje bazy]
\label{thm:lin-charakterystyki-bazy}
Dla dowolnego podzbioru \(B\subseteq V\)
następujące warunki są równoważne:
\begin{enumerate}[label=(\roman*)]
\item \(B\) jest bazą \(V\);
\item \(B\) jest maksymalnym względem zawierania
      liniowo niezależnym podzbiorem \(V\);
\item \(B\) jest minimalnym względem zawierania
      zbiorem rozpinającym \(V\).
\end{enumerate}
\end{twierdzenie}
\begin{proof}
Jeśli \(B\) jest bazą, to żaden wektor spoza \(B\)
nie może zostać dołączony bez powstania zależności,
bo każdy wektor należy do \(\operatorname{span}(B)\).
Zatem baza jest maksymalnie niezależna.
Jeśli \(B\) jest maksymalnie niezależny, a nie rozpina \(V\),
można dobrać \(v\notin\operatorname{span}(B)\)
i powiększyć \(B\) do niezależnego zbioru,
wbrew maksymalności. Daje to (ii)\(\Rightarrow\)(i).

Jeśli \(B\) jest bazą, usunięcie któregokolwiek
\(b\in B\) uniemożliwia rozpięcie \(b\)
przez pozostałe elementy na mocy
Lematu~\ref{lem:lin-niezaleznosc-kryteria}.
Istotnie, gdyby właściwy podzbiór \(C\subsetneq B\)
rozpinał \(V\), to po wybraniu
\(b\in B\setminus C\) również
\(B\setminus\{b\}\supseteq C\) rozpinałby \(V\),
wbrew poprzedniemu wnioskowi.
Zatem baza jest minimalnym zbiorem rozpinającym,
co daje (iii).
Jeśli zbiór \(B\) rozpina \(V\), ale jest zależny,
to dla pewnego \(b\in B\) zachodzi
\(b\in\operatorname{span}(B\setminus\{b\})\).
Po usunięciu \(b\) nadal rozpinamy \(V\),
wbrew minimalności. Otrzymujemy (iii)\(\Rightarrow\)(i).
\end{proof}

\section{Zastosowanie: podprzestrzeń zadana rekurencją}
\label{sec:lin-przyklad-rekurencja}

\begin{przyklad}[Ciągi Fibonacciego z dowolnymi danymi początkowymi]
\label{prz:lin-rekurencja}
Niech \(V=\mathbb R^{\mathbb N_0}\).
Rozważmy zbiór
\[
 W=\{x=(x_n)_{n\geq0}\in V:
                   x_{n+2}=x_{n+1}+x_n
                   \text{ dla każdego }n\geq0\}.
\]
Jest on podprzestrzenią: ciąg zerowy spełnia
rekurencję, a jeśli \(x,y\) ją spełniają,
to dla \(a,b\in\mathbb R\)
\[
 (ax+by)_{n+2}
 =a(x_{n+1}+x_n)+b(y_{n+1}+y_n)
 =(ax+by)_{n+1}+(ax+by)_n .
\]
Zadanie dwóch pierwszych wyrazów
\(x_0=s,\ x_1=t\) wyznacza dokładnie
jeden ciąg z \(W\): następne wyrazy
obliczamy kolejno z rekurencji.
Niech \(u\in W\) ma dane początkowe
\((1,0)\), a \(v\in W\) dane \((0,1)\).
Wtedy każdy \(x\in W\) ma postać
\(x=x_0u+x_1v\): obie strony mają te
same dane początkowe i spełniają
tę samą rekurencję, więc są równe
na mocy indukcji po indeksie \(n\).
Jeśli \(au+bv=0\), porównanie wyrazów
o indeksach \(0,1\) daje \(a=b=0\).
Zatem \((u,v)\) jest bazą \(W\)
i \(\dim_{\mathbb R}W=2\).

Ważny wniosek: \(W\) jest przestrzenią
dwuwymiarową, choć jej wektory są
\emph{nieskończonymi} ciągami.
Liczba współrzędnych używanych do
opisania elementu nie jest sama
w sobie wymiarem przestrzeni.
\end{przyklad}

\section{Podsumowanie zależności logicznych}
\label{sec:lin-podsumowanie}

W rozdziale nie utożsamialiśmy przestrzeni
wektorowej z \(K^n\): istnienie współrzędnych
wynika dopiero z wyboru bazy.
W przypadku skończenie generowanym
udowodniliśmy bez aksjomatu wyboru:
istnienie bazy, lemat o wymianie,
równość liczby elementów baz oraz
jednoznaczność współrzędnych.
Dla dowolnych przestrzeni
istnienie bazy uzyskaliśmy z lematu Zorna.
Dla ustalonego podzbioru \(B\subseteq V\)
otrzymaliśmy równoważne opisy bazy:
\(B\) jest liniowo niezależny i
\(\operatorname{span}(B)=V\) wtedy i tylko wtedy,
gdy każdy \(v\in V\) posiada dokładnie jedno
przedstawienie jako skończona kombinacja liniowa
elementów \(B\).
To ostatnie sformułowanie wymaga
\emph{obu} warunków: istnienia i jednoznaczności
przedstawienia.
W kolejnym rozdziale rozwiniemy pojęcie wymiaru,
sumy i ilorazy podprzestrzeni oraz
konstrukcje nieskończenie wymiarowe.
